Recovering latent environments requires distinguishing mixture-parameter
identification from recording-level assignment recovery. We construct
two partitions, unrelated by label permutation, that disagree on half
the support yet induce the same pooled distribution, demonstrating
non-identifiability without structural restrictions. In a
two-dimensional, equally weighted Gaussian mixture with covariances
symmetric around the identity, fourth-order cumulants identify the
unordered component covariance pair. Separately, we derive a
finite-sample risk-gap guarantee for global empirical \(K\)-means under
fixed \(K\), a pre-given feature transformation, and explicit
boundedness, population-alignment, and regularity assumptions. Alignment
is assumed rather than derived; the guarantee neither establishes
consistent blind \(K\) selection nor controls finite-initialization
optimization error. Using averaged log-covariance features across three
EEG datasets, known-\(K\) clustering on a label-assisted balanced design
of 90 recordings achieved a recording-level adjusted Rand index (ARI) of
\(0.582\pm0.068\) across six prespecified seeds (mean ± standard
deviation). The blind selector chose \(K=2\) and \(K=3\) in three seeds
each, achieving ARI \(0.419\pm0.140\). Unbalanced recording-level and
subject-level plug-in analyses yielded ARIs of \(0.371\) and \(0.241\),
respectively. Equalizing window counts yielded a mean ARI of \(0.577\).
Although metadata alone perfectly predicts source labels, the
fixed-window result indicates that window count alone does not explain
the observed grouping. Remaining acquisition and dataset confounds
restrict interpretation to source-associated technical groups. These
findings distinguish restricted parameter identification from
conditional assignment recovery, while reliable cluster-number selection
and biological or causal interpretation remain unestablished.
